%! Sinusoidal Function Transformations — Unit 3, Lesson 3.6 — Student Packet Cover Sheet
\documentclass[10pt]{article}
\usepackage{precalculus-article}
\usepackage{precalculus-boxes}
\usepackage{ltablex}
\keepXColumns

\begin{document}

% Full-bleed navy banner
\begin{tikzpicture}[remember picture, overlay]
  \fill[navy] (current page.north west)
    rectangle ([yshift=-0.9in]current page.north east);
\end{tikzpicture}
\vspace{-0.8in}

\begin{center}
  {\color{white}\textbf{\LARGE Precalculus}}\\[4pt]
  {\color{skymid}\large\textbf{Unit 3: Trigonometric and Polar Functions}}\\[6pt]
  {\color{white}\textbf{\large Lesson 3.6 \quad Sinusoidal Function Transformations}}
\end{center}

\vspace{0.22in}
\namedateperiod
\vspace{0.18in}

\begin{learningtargetbox}
\small
\begin{enumerate}[leftmargin=1.5em, itemsep=5pt]
  \item I can identify the amplitude, midline (vertical shift), period, and phase shift
    of a sinusoidal function from its equation.
    \hfill\textit{(LO 3.6.A)}
  \item I can explain how each parameter $a$, $b$, $c$, $d$ in $f(\theta)=a\sin(b(\theta+c))+d$
    corresponds to a specific additive or multiplicative transformation of the parent
    sine function. \hfill\textit{(LO 3.6.A)}
  \item I can rewrite a cosine function as a phase-shifted sine function and explain why
    sine and cosine are equivalent up to a horizontal translation.
    \hfill\textit{(LO 3.6.A, EK 3.6.A.1)}
\end{enumerate}
\end{learningtargetbox}

\vspace{0.14in}

\begin{tocbox}
\small
\renewcommand{\arraystretch}{1.5}
\begin{tabularx}{\linewidth}{@{} c l X r @{}}
\rowcolor{sky}
\textbf{\#} & \textbf{Component} & \textbf{Description} & \textbf{Score} \\
\midrule
1 & Warm-Up        & Spiral review: unit-circle values, parent sine/cosine graphs, transformation vocabulary & \blank{1.2cm} \\
2 & Guided Notes   & Parameters $a$, $b$, $c$, $d$; amplitude, period, phase shift, midline; sine--cosine equivalence & \blank{1.2cm} \\
3 & Group Activity & Identify and interpret sinusoidal parameters in context (tiered R/A/E) & \blank{1.2cm} \\
4 & Exit Ticket    & Identify all four parameters from a given equation; match equation to verbal description & \blank{1.2cm} \\
5 & Homework       & Mixed: parameter identification, period/amplitude computation, AP-style MC & \blank{1.2cm} \\
\midrule
  & \multicolumn{2}{r}{\textbf{Total}} & \blank{1.2cm} \\
\end{tabularx}
\end{tocbox}

\end{document}
